\documentclass[10pt,a4paper]{amsart}
\usepackage[margin=19mm]{geometry}
\usepackage{amsmath,amssymb,mathtools}
\usepackage[scr=rsfs,cal=euler]{mathalfa}
\usepackage{xfrac}
\usepackage[unicode,hidelinks,bookmarks=false]{hyperref}
\newcommand{\ie}{i.e.\ }
\newcommand{\cf}{cf.\ }
\newcommand{\resp}{respectively\ }
\newcommand{\Subsection}{\subsection}
\providecommand{\nobreakdash}{\nobreak\hbox{-}}
\setlength{\emergencystretch}{2em}
\multlinegap0pt
\usepackage{bm}
\usepackage{bbm}
\usepackage{dsfont}
\usepackage{xspace}
\usepackage{cancel}
\usepackage{mathtools}
\usepackage{amsthm}
\makeatletter
\@ifundefined{lemma}{\newtheorem{lemma}{Lemma}}{}
\makeatother
\makeatletter
\expandafter\ifx\csname proof\endcsname\relax
\newenvironment{proof}{\noindent{\bf Proof:}}{\endpf}
\fi
\makeatother
\title[Nonparametric inference from possibly unstable M/G/1 workloa]{Nonparametric inference from possibly unstable M/G/1 workload observations}
\author{Royi Jacobovic, Binyamin Kobzantsev}
\date{Source: arXiv:2607.08472v3 (CC BY 4.0)}
\begin{document}
\maketitle

\noindent\emph{Source and licence.} Nonparametric inference from possibly unstable M/G/1 workload observations. Version arXiv:2607.08472v3 (CC BY 4.0), distributed under the Creative Commons Attribution 4.0 International licence, as recorded in the arXiv licence metadata for this exact version and shown on the abstract page https://arxiv.org/abs/2607.08472v3. Full rights notice: licenses/arxiv-2607.08472-CC-BY-4.0.txt.\par\smallskip

\noindent\textit{Editorial scope.} Counted unit: the Lemma whose source label is \texttt{lem: coupling} in the article (source0.tex lines 683--702, source label \texttt{lem: coupling}), delivered together with the article's own complete original proof of it (source0.tex lines 703--897, one \texttt{proof} environment of 195 lines). No mathematics is written, repaired or re-derived: statement and proof are the article's own text line for line. Required context retained uncounted: eq: Z (lines 467--469); eq: W (lines 471--473). Every definition, display and label of those lines is retained unchanged. Omitted from this package and not redistributed: 1--466, 470--470, 474--682, 898--1921. Nothing inside a retained excerpt is omitted or altered: every retained body is delivered exactly as the article's lines are written, so no omission receipt of any kind accompanies this package. Citations: the retained excerpts carry no citation command at all, so no bibliography entry of the article is transcribed and no reference is redistributed. Reference adaptation: none was needed and none was made; every reference inside the delivered unit resolves inside it, so no unresolved reference or citation is produced. Printed slips kept literal: no printed word, symbol, hypothesis or number of the counted statement or of its proof was altered. Layout and class: the article's own class is \texttt{article} with packages including geometry, pgfplots, subfig, amsmath, amssymb, algorithmic, algorithm, mathtools, color, ulem, setspace, tikz, float, hyperref; the wrapper is the lane's pinned \texttt{amsart} editorial wrapper at 10pt on A4, which restates the theorem-like environments the retained text needs and adds the article's own macro definitions that the retained text needs (none), with extra wrapper packages (bm, bbm, dsfont, xspace, cancel, mathtools). The delivered numbering is the unit's own. Assets: no retained span contains a figure environment, a graphics-inclusion command, a PDF-page inclusion, a table or a TikZ picture; no image file is copied into this package, the package declares no asset dependency and no image file is redistributed with it. Licence: version arXiv:2607.08472v3 of this article is distributed under the Creative Commons Attribution 4.0 International licence, as recorded in the arXiv licence metadata for this exact version and shown on the abstract page https://arxiv.org/abs/2607.08472v3. A full rights notice is delivered at licenses/arxiv-2607.08472-CC-BY-4.0.txt. Characters: every retained excerpt is pure ASCII, so the delivered engine, XeLaTeX with its default Latin Modern text font, reports no missing character.
\begin{equation}\label{eq: Z}
Z_t\equiv Z_t\left(J,x\right)\equiv x +J_t-t,\qquad t\geq0,
\end{equation}

\begin{equation}\label{eq: W}
W_t\equiv W_t\left[Z(J,x)\right]\equiv Z_t-\inf_{0\leq s\leq t} Z_s^-,\qquad t\geq0.    
\end{equation}

\begin{lemma}\label{lem: coupling}
There exist  $0<q<p<1$ and three random variables $A_n,K_n^*$ and $V_n$ such that:
\begin{enumerate}
    \item[(i)] $A_n\leq L_n$ almost surely and 
    \begin{equation*}
        A_n\sim\text{\normalfont Bin}(n-m_n,q).
    \end{equation*}

    \item[(ii)] $K_n\leq V_n$ almost surely such that 
    \begin{equation*}
    V_n\sim\text{\normalfont Bin}(m_n+1,p).
    \end{equation*}
    \item[(iii)] $K_n^*\leq K_n$ almost surely such that 
   \begin{equation*}
    K_n^*\sim\text{\normalfont Bin}(\overline m_n+1,p^2).
    \end{equation*}
    
     \item[(iv)] $A_n$ and $(K_n^*,V_n)$ are independent.  \end{enumerate}

\end{lemma}

\begin{proof}
Recall the processes
$t\mapsto Z_t(x,J)$
and
$t\mapsto W_t[Z(J,x)]$
defined in
\eqref{eq: Z}
and
\eqref{eq: W}.
Since
\[
Z_t(x,J)=x+J_t-t,
\]
the mapping
$x\mapsto Z_t(x,J)$
is increasing for every fixed $t\ge0$.
Moreover, the Skorokhod reflection map is monotone with respect to the pointwise order on
$\mathcal D([0,\infty))$.
Consequently,
\[
x\longmapsto W_t[Z(J,x)]
\]
is also increasing.

We now construct an auxiliary workload process by restarting the queue from zero after every observation epoch. Define
\[
W_k^*
=
W_1[Z(J^k,0)],
\qquad
k\ge1,
\]
with $W_0^*=0$.
Since $W_{k-1}\ge0$, monotonicity yields
\[
W_k^*\le W_k,
\qquad
k\ge1.
\]
Hence,
\[
I_k^*
\equiv
\mathbf1_{\{W_k^*>1\}}
\le
I_k,
\qquad
k\ge1,
\]
almost surely.

Observe that $W_k^*$ depends only on $J^k$.
Therefore,
$(I_k^*)_{k\ge1}$
is an i.i.d. sequence.
Furthermore,
\begin{align*}
q
&\equiv
\mathbb P(I_k^*=1)
\\
&=
\mathbb P\!\left(
W_1[Z(J^1,0)]>1
\right)
\\
&>
\mathbb P\!\left(
Z_1(J^1,0)>1
\right)
\\
&=
\mathbb P(J_1^1>2)
\equiv
p.
\end{align*}
In particular, the strict inequality holds since when starting at zero, the reflection remains active until the first jump.
Thus,
\[
I_k^*\sim\mathrm{Ber}(q).
\]

Define
\[
A_n
=
\sum_{k=m_n+1}^{n}I_k^*.
\]
Since $I_k^*\le I_k$,
\[
A_n\le L_n
\]
almost surely, while
\[
A_n\sim\mathrm{Bin}(n-m_n,q),
\]
proving (i).

To establish (ii), observe that
\[
I_k\mathbf1_{\{X_k>2\}}
=
I_k\mathbf1_{\{J^{k+1}_1>2\}}
\le
\mathbf1_{\{J^{k+1}_1>2\}}.
\]
Hence,
\[
V_n\equiv\sum_{k=0}^{m_n}
\mathbf1_{\{J^{k+1}_1>2\}}
\]
satisfies
\[
K_n\le V_n
\]
almost surely.
Since the indicators are i.i.d. Ber$(p)$,
\[
V_n\sim\mathrm{Bin}(m_n+1,p),
\]
which proves (ii).

For (iii), note that
\[
I_k
\mathbf1_{\{X_k>2\}}
=
I_k
\mathbf1_{\{J^{k+1}_1>2\}}
\ge
I_k^*
\mathbf1_{\{J^{k+1}_1>2\}}.
\]
Moreover,
\[
I_k^*
\geq
\mathbf1_{\{J^{k+1}_1>2\}}
,
\]
since
$J^{k+1}_1>2$
implies
$W_k^*>1$.
Therefore,
\[
I_k
\mathbf1_{\{X_k>2\}}
\ge
\mathbf1_{\{J^{k+1}_1>2\}}
\mathbf1_{\{J^{k+1}_1>2\}}.
\]
Define
\[
K_n^*\equiv
\sum_{k=0}^{\overline m_n}
\mathbf1_{\{J^{2k+1}_1>2\}}
\mathbf1_{\{J^{2k+2}_1>2\}}.
\]
The summands involve disjoint pairs of independent processes and are therefore i.i.d. Ber$(p^2)$.
Furthermore,
\[
K_n^*
\le
K_n
\]
almost surely.
Hence,
\[
K_n^*
\sim
\mathrm{Bin}(\overline m_n+1,p^2),
\]
establishing (iii).

Finally,
$A_n$
depends only on
\[
J^{m_n+1},
J^{m_n+2},
\ldots,
\]
whereas
$(K_n^*,V_n)$
depends only on
\[
J^1,\ldots,J^{m_n}.
\]
Since these two collections of driving processes are independent, so are
$A_n$
and
$(K_n^*,V_n)$,
which proves (iv).
\end{proof}

\end{document}
